\uuid{1Vuf}
\exo7id{3504}
\titre{exo7 3504}
\auteur{quercia}
\organisation{exo7}
\datecreate{2010-03-10}
\isIndication{false}
\isCorrection{true}
\chapitre{Réduction d'endomorphisme, polynôme annulateur}
\sousChapitre{Diagonalisation}
\cpge{mpsi-mp2i.nombres-complexes.f, pcsi.nombres-complexes.f, ptsi.nombres-complexes.f, mpsi-mp2i.calcul-matriciel-systemes.e, pcsi.determinants.c, ptsi.determinants.a, mpsi-mp2i.changements-bases.a, pcsi.matrices-applications-lineaires.c, ptsi.matrices-applications-lineaires.c, bcpst1.matrices, mp-mpi.reduction-endomorphismes.b, pc.reduction.a, psi.reduction-endomorphismes.a, pt.reduction-endomorphismes.a, mp-mpi.reduction-endomorphismes.d, pc.reduction.c, psi.reduction-endomorphismes.c, pt.reduction-endomorphismes.c, bcpst2.valeurs-propres-vecteurs-propres.a, bcpst2.valeurs-propres-vecteurs-propres.b}
\module{Algèbre}
\niveau{L2}
\difficulte{}

\contenu{
\question{
Diagonaliser $M = \begin{pmatrix} 0 &       &       &1  \cr
                            1 &\ddots &(0)        \cr
                              &\ddots &\ddots     \cr
                            (0) &     &1      &0  \cr \end{pmatrix} \in \mathcal{M}_n(\C)$.
}
\reponse{
$P = \bigl(\omega^{(i-1)(1-j)}\bigr)$,
         $D = \text{diag}(1,\omega,\dots,\omega^{n-1})$
         avec $\omega = \exp(2i\pi/n)$.
}
}
